\textbf{COLUNI 2019 -- Questão 31}\par\medskip Leia o texto:

\par

\textbf{Tremor de terra de magnitude 3,6 na escala Richter faz Florianópolis tremer}

\par

Os moradores de Florianópolis (SC) sentiram a terra tremer na manhã desta sexta-feira (13/4). Os centros de sismologia das universidades de Brasília (UnB) e de São Paulo (USP) registraram, às 9h28, um tremor de magnitude 3,6 na escala Richter, na margem continental do estado. De acordo com a defesa civil de Santa Catarina, foram relatados tremores no norte e no sul da ilha. Em São João Batista, também chegaram relatos do tremor. “Até o momento a situação está sendo averiguada”, informou o órgão. O prefeito de Florianópolis, Gean Loureiro, disse por meio das redes sociais que, “a princípio, (não há) nenhum problema com estruturas em residências. [...]"

\par

(RESENDE, Humberto. \textbf{Terremoto de magnitude 3,6 na escala Richter faz Florianópolis tremer}. Disponível em: <https://correiobraziliense.com.br/app/moticia/brasil/2018/04/13/interna-brasil,673417/terremoto-demagnitude-3-6-na-escala-richeter-e-sentido-em-florianopolis.html>. Acesso em: 13 abr. 2018. Adaptado.)

\par

A informação sobre o tremor de terra na cidade de Florianópolis (SC) é:\begin{enumerate}
\item falsa, uma vez que no Brasil não existem tremores de terra.
\item procedente, pois no Brasil pode haver tremores de terra.
\item irreal, pois no Brasil os tremores de terra ocorrem apenas na Amazônia.
\item verdadeira, pois Florianópolis está na borda da placa tectônica do Pacífico.\end{enumerate}
